\section{Companion lab: regrouping bytes as Base64 text}
\label{revision-lab:base64}

The delta codec assigns short spellings to small integers. Base64 has a different
purpose: represent arbitrary bytes using a restricted text alphabet. It expands
three input bytes into four characters and uses padding for an incomplete final
group. It neither compresses nor encrypts the data. This companion revises
original project \texttt{028-base64}; use the complete
\texttt{base64-codec.asm} in \texttt{revision/labs/files}. It supplies scalar and
MSA whole-buffer entries plus exact-size vector kernels. The independent
\texttt{base64url.asm} starter changes the alphabet and omits padding.

You will need unsigned byte loads, masks and shifts, function calls, the earlier
MSA lane/shuffle chapter and the distinction between validation and publication.
Read the scalar bit grouping before the vector schedule. These activities use
ordinary Run and Debug; the native guided lesson remains unfinished. Keep source
copies for changed inputs. Original project checks and input recipes remain in
\texttt{packages/028-base64}.

\subsection*{Derive four indices from three bytes}

Let the three unsigned bytes be $a,b,c$. Concatenate their bits in that order,
most significant bit first, then split the 24 bits into four groups of six.
The resulting indices are
\[
 q_0=a\mathbin{\gg}2,\qquad
 q_1=((a\mathbin{\&}3)\mathbin{\ll}4)\mathbin{|}(b\mathbin{\gg}4),
\]
\[
 q_2=((b\mathbin{\&}15)\mathbin{\ll}2)\mathbin{|}(c\mathbin{\gg}6),\qquad
 q_3=c\mathbin{\&}63.
\]
Each $q$ lies between zero and 63. These are numerical indices, often called
\emph{sextets}; they are not yet character codes. The two middle sextets cross
input-byte boundaries, so shifting each byte independently is insufficient.

For bytes hexadecimal \texttt{4d 61 6e}, the text \texttt{Man}, the indices are
$19,22,5,46$. The standard alphabet assigns indices zero through 25 to
\texttt{A} through \texttt{Z}, 26 through 51 to lowercase letters, 52 through
61 to digits, and 62/63 to \texttt{+/}. The four indices become
\texttt{TWFu}, bytes \texttt{54 57 46 75}. The zero index means character
\texttt{A}, ASCII 65, not byte zero. The zero byte remains valid binary input.

\begin{figure}[htbp]
\centering
\begin{tabular}{l l l l}
\toprule
View & First & Second & Third/fourth\\
\midrule
Input bytes & \texttt{01001101} & \texttt{01100001} & \texttt{01101110}\\
Six-bit groups & \texttt{010011} & \texttt{010110} & \texttt{000101 101110}\\
Indices & 19 & 22 & 5, 46\\
Characters & \texttt{T} & \texttt{W} & \texttt{F}, \texttt{u}\\
\bottomrule
\end{tabular}
\caption{Regrouping the same 24 bits of \texttt{Man}. The second and third sextets cross byte boundaries. The alphabet lookup converts each numerical sextet to one character byte; it does not change the six-bit payload.}
\label{fig:revision-base64-groups}
\end{figure}


The scalar encoder loads with \texttt{lbu}. In its first two-character operation,
\texttt{\$t0}=a, \texttt{\$t1}=b, \texttt{\$s4} points to the alphabet and
\texttt{\$s2} points to output. After writing the first character, this fragment
forms and publishes the second:
\Needspace{10\baselineskip}
\MARSCodeLabel{Join two input bytes into sextet one}
\begin{minted}{mars-idea}
andi $t3, $t0, 3
sll $t3, $t3, 4
srl $t4, $t1, 4
or $t3, $t3, $t4
daddu $t4, $s4, $t3
lbu $t4, 0($t4)
sb $t4, 1($s2)
\end{minted}
The masks justify the small word shifts; the alphabet address still uses full-
width arithmetic. A numerical result of 22 becomes ASCII 87 only at the lookup.
Investigate those as separate states if the bit grouping looks correct but the
stored text does not.

Run the unchanged caller, whose mode two selects scalar encoding of
\texttt{Man}. Break at \texttt{encode\_scalar\_return}, before
\texttt{main\_done}. Status \texttt{\$v0}=0 and count \texttt{\$v1}=4 are
live; destination begins \texttt{TWFu}, and its next byte remains the sentinel
\texttt{a5}. This interface writes an explicit length, without a NUL terminator.
The caller publishes status and count at \texttt{main\_done}, then exits.

\subsection*{An incomplete group still has four characters}

One final input byte has eight useful bits. It fills $q_0$ and the high two bits
of $q_1$; the low four bits of $q_1$ must be zero. Emit those two alphabet
characters followed by \texttt{==}. Two final bytes fill $q_0,q_1$ and the high
four bits of $q_2$; its low two bits must be zero. Emit three characters followed
by one \texttt{=}. Padding is a marker, not an extra sextet with value 64.

For \texttt{M}, the result is \texttt{TQ==}; for \texttt{Ma}, it is
\texttt{TWE=}; for \texttt{Man}, \texttt{TWFu}. Predict these before changing
count to one, two or three in a copy. Keep the larger source allocation, but
explain why the encoder must not read missing bytes from that backing storage.
Its scalar path initializes absent $b$ or $c$ to zero and loads them only when
the selected input contains them.

For nonnegative count $n$, the required standard output length is
$4\lceil n/3\rceil$. Three bytes need four output bytes; four need eight.
The whole encoder first rejects an over-limit unsigned count, then accepts
empty input, then checks the complete required capacity before reading source.
Insufficient capacity returns status three with written count zero and no
caller access. A high-bit-set capacity is a large unsigned capacity, not a
negative signed capacity.

\subsection*{Validate all text before publishing decoded bytes}

The decoder accepts canonical standard-alphabet text: length is a multiple of
four, padding appears only in its legal final positions, unused bits are zero,
and whitespace and nonalphabet bytes are rejected. \texttt{TR==} has the same
high useful bits as \texttt{TQ==}, but its second sextet is 17 rather than 16.
Its low four bits are nonzero, so it is rejected. Similarly, \texttt{TWF=}
would discard nonzero low bits from its third sextet and is rejected, while
\texttt{TWE=} is canonical.

The worked helper \texttt{b64\_validate} scans the complete input and computes
the decoded length without writing output. The whole decoder checks capacity
only after validation succeeds, then decodes. Thus malformed text takes precedence
over insufficient capacity; every failure has written count zero and leaves
all output untouched. This is a whole-buffer guarantee, unlike the delta
codec's accepted-prefix guarantee. State which guarantee a caller can rely on.

Prepare a decoder copy with mode three, count eight, source text
\texttt{TWFuTR==} and any sufficient capacity. The first quartet is valid, but
the final quartet has nonzero unused bits. Stop at
\texttt{decode\_scalar\_return}. Expect malformed status two, written count
zero and an entirely unchanged destination. Then set capacity zero and repeat;
malformed still wins. Replace the text with \texttt{TWFuTQ==} and repeat with
capacity zero: now valid text reaches the capacity test and returns status three.

\begin{figure}[htbp]
\centering
\begin{tikzpicture}[alt={Two actual light-theme Debug observations after scalar decoding of TWFuTR double-equals. Registers v0 and v1 show malformed status two and zero written bytes. The displayed first output row at hexadecimal address 10012000 remains sixteen a5 sentinel bytes. A late noncanonical quartet has prevented publication of even the valid first quartet. Native labels and boxes mark the return values and unchanged row.}]
\node[anchor=north west,inner sep=0,font=\small\bfseries] (atitle) at (0,0) {1. Malformed status and zero length};
\node[anchor=north west,inner sep=0] (a) at ([yshift=-2mm]atitle.south west) {\MARSRevisionCrop{revision/captures/base64/base64-invalid-status.png}{0.25\linewidth}{0.24609375}{0.31333333}{0.66015625}{0.64222222}};
\node[anchor=north west,inner sep=0,font=\small\bfseries] (btitle) at ([yshift=-4mm]a.south west) {2. Beginning of the unchanged destination};
\node[anchor=north west,inner sep=0] (b) at ([yshift=-2mm]btitle.south west) {\MARSRevisionCrop{revision/captures/base64/base64-unchanged-output.png}{0.97\linewidth}{0.0328125}{0.3388889}{0.584375}{0.6088889}};
\begin{scope}[shift={(a.south west)},x={($(a.south east)-(a.south west)$)},y={($(a.north west)-(a.south west)$)}]
\draw[black,line width=.7pt] (.035,.02) rectangle (.50,.98);
\end{scope}
\begin{scope}[shift={(b.south west)},x={($(b.south east)-(b.south west)$)},y={($(b.north west)-(b.south west)$)}]
\draw[black,line width=.7pt] (.20,.02) rectangle (.998,.43);
\end{scope}
\end{tikzpicture}
\caption{Actual return-state and Memory captures from the same deterministic late-error input. The first output row remains unchanged; executable access checks separately verify the whole caller buffer and absence of temporary output stores. Native LaTeX supplies labels, crops and boxes.}
\label{fig:revision-base64-native-validation}
\end{figure}


Decoding reverses the bit grouping. For validated sextets, the bytes are
$(q_0\ll2)|(q_1\gg4)$, $((q_1\mathbin{\&}15)\ll4)|(q_2\gg2)$ and
$((q_2\mathbin{\&}3)\ll6)|q_3$. Padding determines which of these bytes
exist; do not store three bytes for every final quartet. Exact in-place decoding
is allowed because each four-character group is read before its shorter output
is stored. Earlier stores cannot reach a future quartet. Shifted partial overlap
is outside this interface, as is in-place encoding that expands into unread input.

\subsection*{Vectorize the same bit argument with exact footprints}

The direct encoder kernel takes exactly twelve source bytes and writes sixteen
characters. Its caller has already established those bytes exist. A sixteen-byte
vector load directly from source would overread by four, even if its unused
lanes were later ignored. The worked kernel copies twelve bytes into a private
sixteen-byte stack slot whose remaining bytes are zero, then performs its vector
load there. The private allocation supports sixteen bytes; the source slice does
not. A different valid schedule may assemble the vector another way.

The shuffle and shift tables construct sixteen sextets. Four groups of three
input bytes become four groups of four output lanes. For \texttt{ManManManMan},
register \texttt{\$w1} at \texttt{after\_enc12\_right} contains byte lanes
$[19,22,5,46]$ repeated four times. The later alphabet conversion produces
\texttt{TWFuTWFuTWFuTWFu} in \texttt{\$w4}. Set mode four and count twelve
in a source copy to observe this direct kernel. Its input/output pointers are
in \texttt{\$a0/\$a1}, unlike the whole entry's \texttt{\$a0/\$a2} pair.

\begin{figure}[htbp]
  \centering
  \begin{tikzpicture}[alt={Actual native eight-bit lane row of w1 for twelve input bytes ManManManMan after vector sextet assembly. Hexadecimal lanes 13,16,05,2e repeat four times, representing numerical indices nineteen, twenty-two, five and forty-six. The native box identifies these indices before conversion to the character bytes TWFu repeated four times.}]
    \node[anchor=south west,inner sep=0] (a) {\MARSRevisionCrop{revision/captures/base64/base64-sextets.png}{0.99\linewidth}{0.2375}{0.262}{0.28125}{0.718}};
    \begin{scope}[shift={(a.south west)},x={($(a.south east)-(a.south west)$)},y={($(a.north west)-(a.south west)$)}]
      \draw[black,line width=.7pt] (.025,.02) rectangle (.985,.98);
    \end{scope}
  \end{tikzpicture}
  \caption{Actual numerical sextets in \texttt{\$w1} before alphabet conversion. Each four-lane group has the same indices as the scalar \texttt{Man} example. Values are shown in hexadecimal; \texttt{13} is decimal nineteen, not ASCII \texttt{T}. Native LaTeX supplies the crop and box.}
  \label{fig:revision-base64-native-sextets}
\end{figure}


\Needspace{8\baselineskip}
The direct decoder starts with exactly sixteen valid, unpadded alphabet
characters and produces twelve bytes. Its vector result initially has unused
positions between three-byte groups. A compaction shuffle gathers the twelve
useful bytes. The kernel stores that vector to private stack storage, then copies
only twelve bytes to the caller. A direct sixteen-byte output store would violate
the interface even when its first twelve bytes were correct. Direct-kernel input
validity is a precondition, not a promise to reject malformed arbitrary text.

Whole MSA entries combine these kernels with the same validation, capacity and
scalar-tail rules as their scalar counterparts. They select the encoder kernel
only when at least twelve raw bytes remain. Decoder vector work applies only to
complete unpadded groups. Explain the boundary at eleven/twelve/thirteen raw
bytes and at the last padded quartet. Reusing a vector operation does not transfer
its stronger preconditions to every whole-buffer input.

\subsection*{Check the complete contract, then compare work}

Whole entries take full source pointer, unsigned count, full destination pointer
and unsigned capacity in \texttt{\$a0} through \texttt{\$a3}. Status in
\texttt{\$v0} is zero success, one count limit, two malformed decoder text or
three capacity. \texttt{\$v1} is exact output length on success and zero on
failure. Encoders accept up to 4096 input bytes, requiring at most 5464 text
bytes. Decoders accept up to 5464 text bytes, which can describe as many as 4098
binary bytes. Those two limits are not inverse bounds on the same set of inputs.

All pointers may be byte-unaligned. Preserve saved registers, frame pointer,
stack pointer and their full bits, return through the supplied continuation and
perform no console I/O. Read selected input/output slices and write only the
exact successful result. Private constants and stack storage are allowed. The
8192-byte source/destination arenas and private workspace are backing capacities,
not permission to overread, overstore or append a terminator.

\Needspace{12\baselineskip}
Run the existing completed-interface suites against the complete source. For
example:
\begin{verbatim}
java -jar /path/to/cli.jar grade \
  packages/028-base64/checks/canonical.kts \
  revision/labs/files/base64-codec.asm
java -jar /path/to/cli.jar grade \
  packages/028-base64/checks/kernels.kts \
  revision/labs/files/base64-codec.asm
\end{verbatim}
The same directory has encoding, decoding, bounds, caller and repeated-session
suites. Use a late error, an exact in-place decode, a maximum input and an empty
input in addition to short familiar words. Check all bytes and preserved state;
a round trip alone cannot establish canonical acceptance or exact footprints.

If comparing scalar and vector work, use the same whole-buffer interval and
include validation, calls, private copies, alphabet conversion and tails. A fast
kernel with extra copying can have a different whole-function cost. Keep
completed instructions separate from host time and hardware cycles, as in the
main chapter's measurement discussion.

\subsection*{Write an unpadded URL-safe encoder independently}

Complete \texttt{b64url\_encode} in \texttt{base64url.asm}, retaining its caller
and allocations. Use the same six-bit groups but map indices 62 and 63 to
\texttt{-\_} instead of \texttt{+/}. Omit all \texttt{=} padding. This is the
URL-safe alphabet and explicitly unpadded variant for this exercise; the name
alone does not imply that padding may always be omitted in another protocol.

Inputs are full source pointer, unsigned byte count zero through 64, full
destination pointer and unsigned capacity in \texttt{\$a0} through
\texttt{\$a3}. Return status zero success, one over-limit count or three capacity
in \texttt{\$v0}, with exact written length or zero in \texttt{\$v1}.
Over-limit count wins before pointer access; empty input then succeeds without
pointer access; capacity is checked before reading nonempty source. Buffers are
disjoint and may begin at arbitrary byte addresses. Preserve other caller data,
saved registers, frame pointer and stack pointer, return through the supplied
continuation and produce no console I/O. No padding, NUL or rejected output may
be written, including temporary stores restored before return.

Derive output length from complete groups and a remainder of zero, one or two:
$4\lfloor n/3\rfloor$ plus zero, two or three. The maximum is 86 bytes for
64 inputs. Do not reserve or report the standard padded length for a short tail.
The default four input bytes are hexadecimal \texttt{fb ff ff ff}, whose
unpadded URL result is \texttt{-\_\_\_\_w}, six bytes. Derive its indices to
check both special symbols and its one-byte remainder.

Establish a scalar implementation, then choose whether and how to reuse vector
sextet generation. Changing a scalar alphabet table does not automatically change
a vector arithmetic classifier. Keep a one-byte input, a two-byte input, all
special-symbol lanes and a twelve-byte group as evidence for your choices.
\Needspace{6\baselineskip}
Run the public checks:
\begin{verbatim}
java -jar /path/to/cli.jar grade \
  revision/labs/checks/base64url.kts \
  /path/to/your/base64url.asm
\end{verbatim}
Checks accept equivalent solutions; they do not require a particular frame,
classifier or vector schedule. Explain the bit grouping, tail length and exact
footprint separately from any performance claim.

The alphabet, grouping and padding conventions follow
\href{https://www.rfc-editor.org/rfc/rfc4648.html}{RFC 4648}; the whole-function
limits, strict decoder policy, precedence and memory guarantees are the lab's
explicit interface. The public package retains its larger construction sequence
and recipes for optional practice. This written companion supplies the reasoning
and independent assignment without requiring that sequence of editor actions.
